\scanpage{124}
\mspage{2}
\deleted{
$Z(S')=$ ensemble des couples $(g,v)\in G(S')\times V(S')$ tels que $v\in g.W_{S'}(S')$.

Par suite, $Z$ est isomorphe au faisceau $G\times_S W$, grâce à l'isomorphisme
\[
G\times_S W\xrightarrow{\sim}Z
\]
défini par $(g,w)\rightsquigarrow(g,g.w)$. Ainsi le diagramme cartésien précédent donne le diagramme cartésien

\diagram{\begin{tikzcd}[ampersand replacement=\&]
\text{(xx)}\&G\times_S W\arrow[r,"q"]\arrow[d,"\lambda"']\&X\arrow[d,"i"]\\
\&G\times_S V\arrow[r]\&G/N\times_S V
\end{tikzcd}}

où on a $\lambda(g,w)=(g,g.w)$, donc $q(g,w)=(g^{\bullet},g.w)$, où $g^{\bullet}$ désigne l'image de $g$ par l'application canonique $G(S')\longrightarrow(G/N)(S')$. On voit enfin sur le diagramme (x) que $Z\longrightarrow X$ fait de $Z$ un fibré principal \deleted{homogène sur} \inserted{de base} $X$ \inserted{et} de groupe $N$ opérant \deleted{\unsure{par}} à droite par $(g,v).n=(gn,v)$, de sorte que \deleted{\unsure{donc}} dans (xx), $q:G\times_S W\longrightarrow X$ fait de $G\times_S W$ un fibré principal de base $X$ et de groupe $N$ opérant \deleted{\unsure{par}} à droite par
\[
(g,w).n=(gn,n^{-1}.w)
\]
Nous \inserted{résumons} \deleted{considérons} \inserted{les principaux morphismes précédents dans} le diagramme suivant :

\diagram{\begin{tikzcd}[ampersand replacement=\&,column sep=large,row sep=large]
\&\&G\times_S W\arrow[d,"q"]\arrow[dr,"\phi"]\&\\
\text{(D)}\&G/N\&X\arrow[l,"p"']\arrow[r,"\psi"']\arrow[d,"i"]\&V\\
\&\&G/N\times_S V\arrow[ul,dashed,"\operatorname{pr}_1"]\arrow[ur,dashed,"\operatorname{pr}_2"']\&
\end{tikzcd}}

où $\psi=\operatorname{pr}_2 i$ et $\phi=\psi q$ \inserted{\begin{tikzpicture}\node[draw,ellipse,inner sep=3pt] {i.e. $\phi(g,w)=g.w$.};\end{tikzpicture}} Si $v$ est une section de $V$ sur $S$, le sous-faisceau $\unsure{X_v}$ de $X$ image inverse de cette section par $\psi$ est donné par

$\unsure{X_v}(S')=$ ensemble des sous-faisceaux $W'$ de $V_{S'}$ qui sont localement conjugués de $W_{S'}$ par $G$, et qui contiennent la section $v_{S'}$ de $V_{S'}$,

tandis que le sous-faisceau $\deleted{\unsure{X_v}}$ de $G\times_S W$ image inverse de $v$ par $\phi$ est \deleted{donné par} isomorphe au sous-faisceau $M_v$ de $G$ \deleted{\unsure{formé}} donné par
}
\margin{$\unsure{M_v}$}
\margin{123}
